% !PRINTED 698
\weberpage{698}
Wenn nun $\gamma = 0$ ist, so folgt hieraus, indem man $t$ und $m$ ins Unendliche wachsen lässt, dass auch $T : t$ den Grenzwerth 0 hat. Ist aber $\gamma$ von 0 verschieden, so ist nach (7)
\[
\operatorname{Lim} \frac{\mu_m}{\mu_{m+1}} = \operatorname{Lim} \frac{m}{m+1} = 1,
\]
und daher nähern sich die beiden Grenzwerthe in (9) dem Werthe $\gamma$, und es folgt also:
\weq{10}{\mathop{\mathrm{Lim}}\limits_{t=\infty} \frac{T}{t} = \gamma.}

Setzen wir zweitens umgekehrt die Grenzgleichung (10) voraus, so wird auch jetzt $\mu_n$ mit $n$ ins Unendliche wachsen müssen, denn sonst würde, da die Gesammtzahl aller $\mu$ unendlich ist, $T$ schon für ein endliches $t$ unendlich, und $\gamma$ könnte nicht endlich sein.

Nehmen wir nun einen Werth $\mu_n$, der in der Reihe der unter einander gleichen
\[
\mu_{m+1}, \; \mu_{m+2}, \; \dots, \; \mu_{m+l}
\]
vorkommt, so dass
\[
m + 1 \leqq n \leqq m + l,
\]
so wird $T$, wenn $t$ durch den Werth $\mu_n$ geht, plötzlich um $l$ Einheiten wachsen, und das Verhältniss $T : t$ wächst um $l : \mu_n$. Da aber $T : t$ einen endlichen Grenzwerth haben soll, so muss
\weq{11}{\operatorname{Lim} \frac{l}{\mu_n} = 0}
sein. Wenn nun $t = \mu_n$ ist, so ist $T = m + l$, und folglich ist
\weq{12}{\frac{T}{t} = \frac{m+l}{\mu_n}, \quad \operatorname{Lim} \frac{m+l}{\mu_n} = \gamma.}

Ferner
\[
\frac{m}{\mu_n} < \frac{n}{\mu_n} \leqq \frac{m+l}{\mu_n},
\]
und folglich ist wegen (11) und (12)
\[
\operatorname{Lim} \frac{n}{\mu_n} = \operatorname{Lim} \frac{m}{\mu_n} = \gamma.
\]
Also ist aus der Gleichung (10) die Gleichung (7) gefolgert \textsuperscript{1}).

Mit Benutzung des Satzes 4. können wir nun dem Satze 3. auch die Form geben:

\begingroup
\footnotesize
\textsuperscript{1}) Dieser Satz ist im Wesentlichen auf dieselbe Weise bewiesen bei \gesperrt{Dirichlet-Dedekind}, Supplement II. Vgl. auch \gesperrt{Riemann's} mathematische Werke, 2. Auflage, Nr. XXX.
\endgroup
